\section[q1]{Question 1}
\subsection{Question 1.a}
\textit{In a single table reporting relevant descriptive statistics in the overall sample
and by year [as a categorical variable, pick two years to perform the analysis].
Include summary statistics as you did in assignment 1. However, you should
include p-values in your table assessing the difference in the following variables
by year.}
\begin{table}[ht]
    \centering
    \begin{tabular}{lcccc}
        \toprule
        Variable & Overall & 2015 & 2020 & P-value \\
        \midrule
        pm10\_concentration & 34.20 ± 36.41 & 37.76 ± 39.33 & 30.78 ± 33.03 & 0.00000 \\
        pm25\_concentration & 13.91 ± 15.25 & 15.11 ± 18.13 & 12.98 ± 12.53 & 0.00925 \\
        no2\_concentration & 15.90 ± 10.61 & 18.16 ± 11.52 & 14.20 ± 9.52 & 0.00000 \\
        pm25\_tempcov & 87.73 ± 22.05 & 86.04 ± 22.99 & 89.17 ± 21.12 & 0.00783 \\
        \bottomrule
\end{tabular}
\caption{Descriptive statistics for WHO regions overall and by year}
\end{table}




\subsection{Question 1.b}
\textit{Write a short paragraph describing any differences you observe in the above
variables across the two years you have selected to analyse}

\begin{itemize}
    \item The mean pm10 concentration decreased from 37.76 in 2015 to 30.78 in 2020. The P-Value shows that we  reject the null hypothesis, indicating that the difference is statistically significant with a p-value of 0.00000.
    \item The mean pm25 concentration also decreased from 15.11 in 2015 to 12.98 in 2020, with a p-value of 0.00925.
    \item The mean no2 concentration showed a notable decrease from 18.16 in 2015 to 14.20 in 2020, with a p-value of 0.00000.
    \item The pm25 temporal coverage increased slightly from 86.04 in 2015 to 89.17 in 2020, with a p-value of 0.00783.
\end{itemize}

\newpage